\def\xlong{20}
\def\ylong{8}
\def\kmax{0.9*\ylong}
\def\ti{0.3*\xlong}
\def\tf{0.65*\xlong}

\begin{tikzpicture}[
    texto/.style={black, font=\sffamily},
    lim ancho/.style={text width=2.5cm, align=center},
    eje/.style={thick, arrows=-Stealth},
    limites/.style={dashed, DarkRed},
    marca/.style={loosely dotted, DarkRed},
    linea central/.style={ultra thick, DarkBlue}
]

\draw[eje] (0,0) -- (\xlong,0) node[midway, below=5pt, texto] {Tiempo ($t$)};
\draw[eje] (0,0) -- (0,\ylong) node[midway, rotate=90, above, texto] {Temperatura ($T$)};

\draw[linea central] (0,0.1*\kmax) to [out=60,in=190] (\ti,\kmax) -- (\tf,\kmax) to[out=-80,in=160] (\xlong,0.1*\kmax);

\draw[limites] (\ti,0) -- (\ti,1.1*\kmax);
\draw[limites] (\tf,0) -- (\tf,1.1*\kmax);

\node[texto, lim ancho] at (0.4*\ti, 0.5*\kmax) {Fase de Calentamiento};
\node[texto, lim ancho] at (0.5*\xlong, 0.6*\kmax) {Fase de Sostenimiento};
\node[texto, lim ancho] at (\tf+0.4*\tf, 0.45*\kmax) {Fase de Enfriamiento};

\end{tikzpicture}